\documentclass[12pt]{article}
\usepackage[a5paper,margin=2cm]{geometry}
\usepackage{lmodern}
\usepackage[T1]{fontenc}
\usepackage{amsmath}
\usepackage[style=numeric,backend=biber]{biblatex}
\usepackage{hyperref}
\addbibresource{refs.bib}

\title{Lantern: how builds scale}
\author{A. Learner}
\date{}

\begin{document}
\maketitle

\section{Introduction}
\label{sec:introduction}

Lantern turns a folder of Markdown notes into a website. Every note
becomes a page, every heading an anchor, and every link between notes
a link between pages. The tool is small on purpose: one command, no
configuration and no state of its own.

This paper asks a single question: how does the time of a build grow
with the number of notes? A tool that people run after every edit has
to stay fast as their notes pile up, so the answer decides whether
Lantern needs a cache.

\section{Method}
\label{sec:method}

We built three folders of notes, with one hundred, one thousand and
ten thousand notes, five times each, and kept the median time of each
folder. The machine was otherwise idle, and the files were read once
before the first build so that every run found them in memory.

Each note was a copy of the same short page: a heading, three
paragraphs, a list and two links to other notes. Copies keep the
measurement honest, because a larger folder then means more work of
exactly the same kind and nothing else.

\section{Measurements}
\label{sec:measurements}

Table~\ref{tab:times} lists the median build time of each folder. Ten
times as many notes take ten times as long, give or take the start-up
cost, which is the same for every folder.

\begin{table}[h]
  \centering
  \begin{tabular}{rr}
    \hline
    Notes & Seconds \\
    \hline
    100 & 0.2 \\
    1000 & 1.9 \\
    10000 & 19.4 \\
    \hline
  \end{tabular}
  \caption{Median build time by number of notes.}
  \label{tab:times}
\end{table}

The spread between the five runs of a folder was small, under three
percent of the median in every case, so the medians can be compared
directly.

\section{Discussion}
\label{sec:discussion}

Parsing dominates the cost of a build. Writing the pages is cheap,
and copying the images takes no measurable time at all. A cache of
parsed notes would therefore remove most of the work of a second
build, at the price of a directory of state that Lantern has so far
managed without.

The linear growth is good news for people with a few hundred notes,
who will never notice the build. People with tens of thousands of
notes will, and for them a cache is worth its complexity.

A last observation concerns links. Checking that every link points to
an existing note needs the whole folder at once, so it cannot be
cached per note. It is, however, a simple lookup in a table of names,
and it stayed below one percent of the build in every run.

\section{Conclusion}
\label{sec:conclusion}

Builds grow linearly, parsing is the cost, and a cache of parsed notes
is the obvious next step. Lantern stays small until someone needs it.

\printbibliography
\end{document}
